\section{求解后的四项核验}\label{ic:verification}
\begin{method}{从答案反查全部步骤}
\textbf{第一项：微分。}把回代后的$F(x)$逐项求导，化成同一分母，验证$F'=f$；只做数值抽查不能替代恒等式。
\par\textbf{第二项：区间。}检查分母、对数和根式限制，记录在哪个连通区间使用同一个常数。
\par\textbf{第三项：分支。}反三角、反双曲和特殊函数均使用已声明分支；把两式合并时检查是否只在局部相差常数。
\par\textbf{第四项：退化。}检查参数为0、重根和换元分母为0的情况，回到未约化的表达式重新分类。
\end{method}
\begin{warning}疑难点：原函数表达式看似不同，可能只差常数；在连通区间上验证两者导数相同，再选一个点确定差值。不能据此把经过原奇点的不同区间强行合并。\end{warning}
